% Draft subsection for the Results section of main.tex.
% Suggested placement: after the 3.2 subsection
% (\subsection{Synthetic images reproduce cohort-level geometric phenotypes}).
%
% Source of numbers: results/Privacy Analysis Summary.pptx (2026-07-14) and the
% face_audit package in results/package.zip. The raw result files
% (nnaa_embeddings.json, nnaa_lpips.json, below_threshold_*.json,
% triplet_accuracy_*.json, calibration_*.json) are NOT in this repository, so
% every value below traces to either the FAR table on slide 15 (exact) or to a
% plot axis (read off, marked inline).
%
% TODO before merging into main.tex:
%   - replace every value marked "% READ FROM PLOT" with the exact figure from
%     nnaa_embeddings.json / nnaa_lpips.json / below_threshold_lpips.json
%   - state the bootstrap replicate count in the sentence that reports the
%     confidence intervals (config default of 10 is a smoke-test value)
%   - RESOLVE the pair-count inconsistency: the ArcFace KDE reports 93
%     same-person pairs, the strict ordinal check reports 186. Likely ordered
%     versus unordered pairs, or two different age-filter settings. One
%     convention must be used in both places
%   - fill <UNVERIFIED> patient counts for the training and held-out partitions
%   - confirm 999 = 10 cohorts x 100 sampled images minus one detection failure,
%     and report the detection-failure count for the real partitions too
%   - figures to be regenerated as vector PDF from the face_audit outputs:
%       kde_density_arcface.png, kde_density_lpips.png,
%       below_threshold_multi_backbone.png, nnaa_embeddings.png, nnaa_lpips.png
%     Do NOT reuse the overlap-band example panel from slide 7: it shows
%     identifiable patient photographs
%   - Discussion should carry the sentence that this is an empirical audit and
%     not a formal privacy guarantee

\subsection{Synthetic images do not reproduce training identities}

Because the generators are distributed with FaceKit, we finally asked whether
the images they produce disclose the identity of the patients whose photographs
trained them. For the ten rare-disease cohorts we compared 999 sampled synthetic
images against the training partition of each cohort, using as a reference the
611 held-out real images that belong to the same cohorts but were excluded from
generator training. Face recognition embeddings separate these cohorts less
cleanly than they separate adults in the benchmarks they were calibrated on:
across 93 same-person and 111,598 different-person pairs, the two ArcFace
distance distributions were largely distinct but retained a tail overlap of
0.0716 (Fig.~\ref{fig:privacy}\subref{fig:privacy-kde-arcface}), so no threshold
separates identities without error. We therefore did not fix a verification
threshold. Each threshold was instead read off the held-out distances
themselves, so that an operating point is defined by the false acceptance rate
it produces among real images of the same cohort, and the synthetic images were
asked only to stay at or below that rate. A held-out image is subject to the
same recognition errors and to the same coincidental resemblance between
children who share a syndrome as a synthetic image is, so what remains above
this reference is the excess attributable to the generator. The comparison
requires the embedding to rank identities correctly rather than to be
calibrated, and it does: for 182 of 186 same-person pairs (97.8\%) the ArcFace
distance was smaller than every different-person distance involving those
images.

Against this reference the synthetic images carried no excess identity signal
(Table~\ref{tab:privacy-far},
Fig.~\ref{fig:privacy}\subref{fig:privacy-far}). At an operating point
corresponding to a 1\% false acceptance rate among real images, 0.10\% of
synthetic images fell below the ArcFace threshold, 0.10\% below the AdaFace
threshold and none below the LVFace threshold. The margin persisted as the
operating point was relaxed: at 5\% the three rates were 1.30\%, 0.70\% and
1.50\%, and at 20\% they were 15.02\%, 11.61\% and 6.71\%. No backbone exceeded
its reference rate at any of the five operating points, and the ordering of the
three models was stable throughout. Within the resolution of these three
recognition models, and at every operating point examined, sampling from the
released generators is therefore not more likely to produce a face matching a
training patient than drawing another real photograph of the same syndrome is.

The picture on the appearance axis is different, and the two must not be
conflated. LPIPS distances between images of one patient and between images of
different patients are almost indistinguishable, overlapping by 0.8438 across
897 same-person and 78,388 different-person pairs
(Fig.~\ref{fig:privacy}\subref{fig:privacy-kde-lpips}); the metric responds to
texture, pose and illumination rather than to identity. Measured with it, the
synthetic images were systematically closer to the training photographs than the
held-out images were, and by a wide margin: 14.1\% of synthetic images fell
below the threshold set at a 5\% reference rate, and 36.1\% below the threshold
set at 20\%. % READ FROM PLOT (below_threshold_lpips)
A generator trained on a cohort reproduces the acquisition characteristics of
that cohort, and the synthetic images inherit them. Because the same metric
cannot tell two photographs of one patient from photographs of two different
patients, this proximity describes the distribution the images were drawn from
and not the individuals within it.

Both axes are summarized by the nearest neighbour adversarial accuracy
(Fig.~\ref{fig:privacy}\subref{fig:privacy-nnaa}). Computed on embeddings, the
privacy loss was $-0.000$ for ArcFace, $-0.002$ for AdaFace and $-0.006$ for
LVFace, with bootstrap confidence intervals that included zero in all three
cases; % READ FROM PLOT (nnaa_embeddings)
synthetic images lie no closer to the training partition than to the held-out
partition. Computed on LPIPS distances the privacy loss was $+0.046$ with a
confidence interval of $[0.031, 0.093]$ that excludes zero, % READ FROM PLOT (nnaa_lpips)
reproducing in a single number the divergence between the two axes. One feature
of these values deserves comment: the adversarial accuracies themselves ranged
from 0.947 to 0.980 on embeddings, % READ FROM PLOT (nnaa_embeddings)
far above the value of one half that would indicate two indistinguishable sets,
because synthetic faces are globally separable from real ones in recognition
embedding space. The absolute level therefore carries no privacy
interpretation here, and only the difference between the two comparisons does.

\begin{table}[t]
  \centering
  \caption{\textbf{Synthetic images flagged against a real-image reference
  rate.} Each row fixes an operating point by taking the threshold from the
  $p$-th percentile of the distances between held-out real images and the
  training partition, so that $p$ is the false acceptance rate the threshold
  produces among real images of the same cohort; entries give the percentage of
  the 999 synthetic images falling below that threshold. Thresholds are cosine
  distances for the three recognition models and LPIPS distances for the
  perceptual metric. Values at or below $p$ indicate no excess over the
  reference.}
  \label{tab:privacy-far}
  \begin{tabular}{cccccccc}
    \toprule
    & \multicolumn{3}{c}{Identity (\% flagged)} & \multicolumn{3}{c}{Threshold} & Appearance \\
    \cmidrule(lr){2-4} \cmidrule(lr){5-7} \cmidrule(lr){8-8}
    $p$ (\%) & ArcFace & AdaFace & LVFace & ArcFace & AdaFace & LVFace & LPIPS (\% flagged) \\
    \midrule
     1 &  0.10 &  0.10 & 0.00 & 0.3758 & 0.3708 & 0.5146 &  3.8 \\
     2 &  0.50 &  0.20 & 0.00 & 0.4115 & 0.3879 & 0.5450 &  7.7 \\
     5 &  1.30 &  0.70 & 1.50 & 0.4386 & 0.4411 & 0.5946 & 14.1 \\
    10 &  5.01 &  4.00 & 2.60 & 0.4750 & 0.5027 & 0.6249 & 23.0 \\
    20 & 15.02 & 11.61 & 6.71 & 0.5274 & 0.5504 & 0.6673 & 36.1 \\
    \bottomrule
  \end{tabular}
\end{table}

% Figure skeleton. Panels (a)-(d) are regenerated from the face_audit outputs
% listed in the TODO block above; \subref requires subcaption.
\begin{figure}[t]
  \centering
  \includegraphics[width=\linewidth]{image/privacy_assessment.pdf}
  \caption{\textbf{Privacy assessment of the released generators.}
  \protect\subref{fig:privacy-kde-arcface} Distribution of ArcFace cosine
  distances between images of one patient and between images of different
  patients, with the overlap between the two densities shaded.
  \protect\subref{fig:privacy-kde-lpips} The same distributions measured with
  LPIPS, which overlap almost completely and therefore carry little identity
  information. \protect\subref{fig:privacy-far} Percentage of synthetic images
  falling below a threshold set at the $p$-th percentile of the held-out
  distances, for three recognition models; the dashed line is the rate the same
  threshold produces among real held-out images.
  \protect\subref{fig:privacy-nnaa} Privacy loss, the difference in nearest
  neighbour adversarial accuracy between the held-out and the training
  comparison, with bootstrap confidence intervals; values above zero indicate
  synthetic images lying closer to the training data than held-out real images
  do.}
  \label{fig:privacy}
\end{figure}
